% ============================================================
%  TEMA 3 — APARTADO 4: CRITERIOS PARA EVALUAR LA PLANIFICACIÓN
%                       Y SISTEMAS DE INFORMACIÓN DE RRHH (SIRH)
%
%  RECORDATORIO AL EQUIPO:
%
%  CRITERIOS DE EVALUACIÓN DE LA PLANIFICACIÓN DE RRHH:
%  La planificación de RRHH se evalúa mediante indicadores (KPIs) que miden
%  su eficacia y eficiencia. Los más habituales son:
%
%    Indicadores de flujo de plantilla:
%      - Tasa de rotación (turnover rate): % de bajas sobre plantilla media.
%      - Tasa de absentismo: horas no trabajadas / horas contratadas.
%      - Tiempo medio de cobertura de vacantes (time-to-fill).
%      - Tasa de promoción interna.
%
%    Indicadores de calidad del ajuste:
%      - Porcentaje de posiciones clave cubiertas por planes de sucesión.
%      - Tasa de retención en el primer año.
%      - Tasa de superación del período de prueba.
%
%    Indicadores de coste:
%      - Coste medio de contratación (cost-per-hire).
%      - Coste del absentismo.
%      - Retorno de la inversión en formación (ROI de formación).
%
%    Indicadores de previsión:
%      - Desviación entre plantilla planificada y real.
%      - % de posiciones cubiertas internamente vs. externamente.
%
%  SISTEMAS DE INFORMACIÓN DE RRHH (SIRH / HRIS):
%  Un SIRH es una plataforma de software que integra y gestiona los datos
%  de los empleados y los procesos de RRHH. Sus módulos típicos son:
%    - Administración de personal (contratos, nóminas, seguridad social).
%    - Gestión del tiempo y asistencia.
%    - Selección y reclutamiento (ATS - Applicant Tracking System).
%    - Formación y desarrollo (LMS - Learning Management System).
%    - Evaluación del desempeño.
%    - Planificación de plantilla y People Analytics.
%
%  Los sistemas más conocidos del mercado:
%    SAP SuccessFactors, Workday, Oracle HCM, Cornerstone OnDemand,
%    Meta4 / NGA HR (muy usado en España), A3 RRHH.
%
%  PARA ESTA SECCIÓN INDICAD:
%  - ¿Qué KPIs usa CaixaBank para medir la planificación de RRHH?
%  - ¿Qué sistema(s) SIRH utiliza y qué módulos tiene activos?
%  - ¿Cómo se podría mejorar el sistema o los criterios de evaluación?
% ============================================================

\section{Criterios de evaluación de la planificación y sistemas de información de RRHH}

\subsection{Criterios e indicadores para evaluar la planificación de RRHH}

La eficacia de la planificación de recursos humanos se mide a través de
un conjunto de indicadores clave (KPIs) que permiten detectar desviaciones
y tomar decisiones correctoras.

\subsubsection{KPIs utilizados por CaixaBank}

\begin{table}[H]
  \centering
  \small
  \caption{Indicadores de planificación de RRHH en CaixaBank}
  \begin{tabularx}{\textwidth}{p{0.35\textwidth} p{0.18\textwidth} X}
    \toprule
    \textbf{Indicador}                    & \textbf{¿Se mide?} & \textbf{Valor aproximado / fuente} \\
    \midrule
    Tasa de rotación (turnover)           & [Sí/No] & [\% — fuente: Memoria Anual / entrevista] \\
    Tasa de absentismo                    & [Sí/No] & [\%] \\
    Tiempo medio de cobertura (TTF)       & [Sí/No] & [días] \\
    Tasa de promoción interna             & [Sí/No] & [\%] \\
    Coste medio de contratación           & [Sí/No] & [\euro{}] \\
    Desviación plantilla real vs. plan    & [Sí/No] & [\%] \\
    Cobertura de sucesión (puestos clave) & [Sí/No] & [\%] \\
    ROI de formación                      & [Sí/No] & [\%] \\
    \bottomrule
  \end{tabularx}
\end{table}

[Análisis de los indicadores más relevantes obtenidos en la entrevista.
Para algunos datos públicos, consultad la Memoria Anual de CaixaBank.]

\subsubsection{Propuestas de mejora en los criterios de evaluación}

% TODO: ¿Qué KPIs deberían medir y no miden? Por ejemplo:
%   - NPS de empleados (Employee Net Promoter Score).
%   - Índice de diversidad e inclusión (si no lo tienen).
%   - Indicadores de bienestar (wellbeing index).
%   - Métricas de People Analytics predictivo.

\begin{itemize}
  \item [Propuesta 1 — justificación]
  \item [Propuesta 2 — justificación]
\end{itemize}

\subsection{Sistemas de Información de Recursos Humanos (SIRH)}

\subsubsection{¿Qué sistema SIRH utiliza CaixaBank?}

% TODO: CaixaBank, dada su escala, es probable que use uno de los grandes
% sistemas enterprise (SAP SuccessFactors, Workday, Oracle HCM).
% En 2021, tras la fusión, habrán tenido que integrar los sistemas de
% CaixaBank y Bankia. ¿Cuál es el sistema unificado?

[Descripción a completar con la información de la entrevista]

\begin{table}[H]
  \centering
  \small
  \caption{Módulos del SIRH de CaixaBank}
  \begin{tabularx}{\textwidth}{p{0.40\textwidth} p{0.12\textwidth} X}
    \toprule
    \textbf{Módulo / Funcionalidad}          & \textbf{¿Activo?} & \textbf{Observaciones} \\
    \midrule
    Administración de personal / nómina       & [Sí/No] & \\
    Gestión del tiempo y asistencia           & [Sí/No] & \\
    Reclutamiento y selección (ATS)           & [Sí/No] & \\
    Gestión del desempeño                     & [Sí/No] & \\
    Formación y desarrollo (LMS)              & [Sí/No] & \\
    Planificación de plantilla                & [Sí/No] & \\
    People Analytics e informes               & [Sí/No] & \\
    Gestión de la compensación                & [Sí/No] & \\
    Portal del empleado (autoservicio)        & [Sí/No] & \\
    \bottomrule
  \end{tabularx}
\end{table}

\subsubsection{¿Cómo mejorar el SIRH?}

% TODO: Una vez descripto el sistema actual, reflexionad sobre qué capacidades
% adicionales aportarían valor. Por ejemplo:
%   - Integrar People Analytics predictivo (predecir quién va a marcharse).
%   - Mejorar la experiencia del empleado en el portal de autoservicio.
%   - Integrar IA para el cribado de CV en selección (pero con control del sesgo).
%   - Mejorar la interoperabilidad entre el SIRH y otros sistemas del banco.

[Descripción a completar]

\begin{notaequipo}
  Preguntad en la entrevista:
  \begin{enumerate}
    \item ¿Qué sistema SIRH utilizan? ¿Ha cambiado tras la fusión con Bankia?
    \item ¿Qué módulos tienen activos? ¿Qué módulos les gustaría activar?
    \item ¿Qué KPIs de plantilla monitorizan de forma habitual?
    \item ¿Tienen un cuadro de mando (dashboard) de RRHH? ¿Quién lo consulta?
    \item ¿Están satisfechos con la información que les proporciona el SIRH?
          ¿Qué mejorarían?
  \end{enumerate}
\end{notaequipo}
